\documentclass[11pt]{article}
\usepackage[margin=1.1in]{geometry}
\usepackage{amsmath,amssymb,amsthm}
\usepackage{booktabs}
\usepackage{microtype}
\usepackage[colorlinks=true,linkcolor=blue,citecolor=blue]{hyperref}

\newtheorem{proposition}{Proposition}
\newtheorem{lemma}{Lemma}
\theoremstyle{remark}
\newtheorem{remark}{Remark}

\newcommand{\dd}{\,\mathrm{d}}
\newcommand{\Om}{\Omega}
\newcommand{\pd}[2]{\frac{\partial #1}{\partial #2}}
\newcommand{\R}{\mathbb{R}}

\title{Corner-Adapted Quadrature for Rellich Normalization\\
\large Why the current rule fails at reentrant corners, and what replaces it}
\author{lappy development notes}
\date{}

\begin{document}
\maketitle

\begin{abstract}
The Rellich identity computes $\|u\|_{L^2(\Om)}$ from boundary Cauchy data alone.
Its integrand carries a factor $r^{2\nu-2}$ at a corner of interior angle
$\alpha$, where $\nu=\pi/\alpha$; at a reentrant corner $\nu<1$ and this is
singular. We show (i) why the singularity can be removed \emph{geometrically} by
placing the Rellich reference point $x_0$ at the corner, (ii) why the current
graded rule cannot integrate it when $x_0$ is placed elsewhere --- the failure is
truncation, not under-resolution --- and (iii) that a corner-local Gauss--Jacobi
rule in the variable $t=r^{\sigma}$ integrates it to machine precision with $O(10)$
nodes when $\sigma$ rationalizes the corner's exponent family, and at a high-order
algebraic rate otherwise --- removing the constraint on $x_0$ in either case.

\medskip\noindent
\textbf{Read Section~\ref{sec:addendum} before implementing.} Sections 1--6 are
the original analysis, and its central claim survived; three of its subsidiary
conclusions did not. The substitution that rationalizes the exponent family is
$t=r^{1/q}$ rather than $t=r^{\nu}$, and is unusable in that form for $q\ge4$; the
$2/\nu\in\mathbb{Z}$ condition is a limitation of $t=r^{\nu}$ specifically, not of
the method; and the recipe of Section~6 omits a geometric cap on panel length whose
absence costs twelve orders on some domains.
\end{abstract}

\section{The identity and its corner integrand}

For $-\Delta u=\lambda u$ in $\Om\subset\R^2$ with $u=0$ on $\partial\Om$, the
Rellich--Pohozaev identity with reference point $x_0\in\mathbb{C}$ gives the
$L^2$ norm as a pure boundary integral,
\begin{equation}\label{eq:rellich}
  \|u\|_{L^2(\Om)}^2
  \;=\; \frac{1}{2\lambda}\int_{\partial\Om}
        \big((x-x_0)\cdot n\big)\Big(\pd{u}{n}\Big)^{\!2}\dd s
  \;=\;\frac{1}{2\lambda}\int_{\partial\Om} r_N \Big(\pd{u}{n}\Big)^{\!2}\dd s ,
\end{equation}
with $r_N:=(x-x_0)\cdot n$. The identity is exact for \emph{every} $x_0$; the
choice of $x_0$ affects only how hard the integral is to evaluate.

\subsection{Local structure at a corner}

Let the corner sit at the origin with interior angle $\alpha$ and edges along
$\theta=0$ and $\theta=\alpha$, and set
\begin{equation}
  \nu \;:=\; \frac{\pi}{\alpha}, \qquad \nu_k := k\nu = \frac{k\pi}{\alpha},
  \quad k=1,2,\dots
\end{equation}
The corner is \emph{reentrant} when $\alpha>\pi$, i.e.\ $\nu<1$; the
$270^\circ$ corner has $\nu=2/3$ and the slit limit $\alpha\to2\pi$ gives
$\nu\to1/2$. Any Dirichlet solution admits the local expansion
\begin{equation}\label{eq:expansion}
  u(r,\theta) \;=\; \sum_{k\ge1} c_k\, J_{\nu_k}\!\big(\sqrt{\lambda}\,r\big)\,
                    \sin(\nu_k\theta).
\end{equation}
On the edge $\theta=0$ the outward normal is $n=-e_\theta$, so $u\equiv0$ there
forces the normal derivative to come entirely from the angular derivative:
\begin{equation}\label{eq:dudn}
  \pd{u}{n}\bigg|_{\theta=0}
  = -\frac1r\pd{u}{\theta}\bigg|_{\theta=0}
  = -\sum_{k\ge1} c_k\,\nu_k\,\frac{J_{\nu_k}(\sqrt\lambda r)}{r}.
\end{equation}
Using $J_{\mu}(z)=(z/2)^{\mu}A_\mu(z^2)$ with $A_\mu$ entire and $A_\mu(0)\neq0$,
\begin{equation}\label{eq:dudn2}
  \pd{u}{n}\bigg|_{\theta=0} = r^{\nu-1}F(r),
  \qquad
  F(r)=\sum_{k\ge1}b_k\, r^{(k-1)\nu}A_{\nu_k}(\lambda r^2),
\end{equation}
so that
\begin{equation}\label{eq:G}
  \Big(\pd{u}{n}\Big)^{\!2}\bigg|_{\theta=0} = r^{2\nu-2}\,G(r),
  \qquad
  G(r) \;=\; F(r)^2 \;=\; \sum_{p\ge0}\sum_{q\ge0} g_{pq}\, r^{\,p\nu+2q}.
\end{equation}
Two families of exponents appear in $G$: the \emph{corner} exponents $r^{p\nu}$
generated by the expansion \eqref{eq:expansion}, and the \emph{Bessel} exponents
$r^{2q}$ from the entire factors. This double structure is the crux of
Section~\ref{sec:sub}.

\subsection{The geometric escape: $r_N\equiv0$}

\begin{proposition}\label{prop:rn}
If $x_0$ is placed at the corner, then $r_N\equiv0$ on both straight edges
emanating from it.
\end{proposition}
\begin{proof}
On the edge $\theta=0$, $x-x_0=r e^{i\cdot 0}$ is parallel to the edge, while
$n=-e_\theta$ is perpendicular to it; hence $(x-x_0)\cdot n=0$ pointwise, and
likewise on $\theta=\alpha$.
\end{proof}

\noindent
The singular factor $r^{2\nu-2}$ is therefore multiplied by an \emph{exact}
zero, and the corner contributes nothing: the boundary integral collapses onto
the parts of $\partial\Om$ where $u$ is smooth. This is why placing $x_0$ at a
reentrant corner yields machine precision with any reasonable quadrature. It is
also why the escape does not generalize: $x_0$ is a single point, and a domain
with several reentrant corners (H-shape, GWW) cannot zero $r_N$ at all of them.

\section{Why the graded rule fails for $x_0$ off the corner}

For general $x_0$, $r_N$ is a nonzero constant along each straight edge (it is
the signed distance from $x_0$ to the edge's line), and the edge contributes
\begin{equation}\label{eq:edge}
  I_{\text{edge}} \;=\; \rho\int_0^R r^{2\nu-2}G(r)\dd r,
  \qquad \rho := r_N = \text{const}.
\end{equation}
This converges for $\nu>1/2$, but the convergence is deceptive.

\begin{lemma}[Mass concentration]\label{lem:mass}
The fraction of $\int_0^R r^{2\nu-2}\dd r$ contributed by $r<\varepsilon$ is
$(\varepsilon/R)^{2\nu-1}$.
\end{lemma}

\noindent
The exponent $2\nu-1$ vanishes as $\nu\to1/2$. Numerically, with $R=1$ and
$\varepsilon=10^{-9}$:
\medskip

\begin{center}
\begin{tabular}{lccc}
\toprule
$\alpha$ & $\nu$ & $2\nu-1$ & mass below $10^{-9}$\\
\midrule
$1.25\pi$ & $0.800$ & $0.600$ & $4.0\times10^{-6}$\\
$1.5\pi$  & $0.667$ & $0.333$ & $1.0\times10^{-3}$\\
$1.75\pi$ & $0.571$ & $0.143$ & $5.2\times10^{-2}$\\
$1.9\pi$  & $0.526$ & $0.053$ & $0.336$\\
$1.99\pi$ & $0.503$ & $0.005$ & $0.901$\\
\bottomrule
\end{tabular}
\end{center}
\medskip

\noindent
For a near-slit corner, \emph{ninety percent of the integral lives within
$10^{-9}$ of the corner.} lappy's Kress-graded rule clamps its innermost node at
\verb|quad._KRESS_TAU_FLOOR| $=10^{-9}$ (a guard against a genuine defect: nodes
built as $p_0+\tau L\,\hat{d}$ in absolute coordinates round to exactly $p_0$
once $\tau L$ falls below machine epsilon relative to $|p_0|$). The rule is thus
\emph{truncated}, not under-resolved: refining from $565$ to $6790$ nodes leaves
the innermost node at $10^{-9}$ and the error unmoved. Observed errors track
Lemma~\ref{lem:mass} over eight orders of magnitude, scaling linearly in
$|x_0-\text{apex}|$ with constant $\approx(10^{-9})^{2\nu-1}$.

\begin{remark}
The floor is a symptom of \emph{coordinate representation}, not of floating-point
resolution. The value $r=10^{-30}$ is perfectly representable; it is
$p_0 + \tau L\hat d$ that is not. A rule evaluated in corner-local $(r,\theta)$
never forms that difference. This distinction is what makes
Section~\ref{sec:jacobi} possible.
\end{remark}

\section{A corner-local Gauss--Jacobi rule}\label{sec:jacobi}

Gauss--Jacobi quadrature with weight $(1-x)^{a}(1+x)^{b}$ integrates
$\int_{-1}^{1}(1-x)^a(1+x)^b f(x)\dd x$ exactly for $f$ polynomial of degree
$\le2n-1$, and spectrally for $f$ analytic. Mapping $r=R(1+x)/2$ and matching
$b=\beta:=2\nu-2$ to \eqref{eq:edge},
\begin{equation}\label{eq:jacobi}
  \int_0^R r^{2\nu-2}G(r)\dd r
  \;=\;\Big(\frac R2\Big)^{\!2\nu-1}\sum_{i=1}^{n} w_i\, G(r_i),
  \qquad r_i=\tfrac R2(1+x_i).
\end{equation}
Admissibility requires $\beta>-1$, i.e.\ $\nu>1/2$: precisely the range of
integrability, degenerating exactly at the slit.

Two properties matter as much as the accuracy:
\begin{enumerate}
\item \textbf{The nodes stay away from the corner.} The smallest Jacobi node
sits at distance $O(n^{-2})$, not $10^{-9}$: for $n=16$ the innermost node is
near $10^{-3}$. The singularity is absorbed \emph{analytically} into the weight
rather than resolved by clustering, so neither the truncation floor nor the
coordinate cancellation of Remark~1 ever arises.
\item \textbf{Only $\nu$ is needed.} $G(r)=(\partial_n u)^2 r^{2-2\nu}$ is formed
numerically from black-box Cauchy data; no analytic knowledge of $u$ is used.
Since $\nu=\pi/\alpha$ comes exactly from the geometry, this is free --- but it
must be exact, see Section~\ref{sec:sens}.
\end{enumerate}

\subsection{The multi-exponent obstruction and the substitution}\label{sec:sub}

Rule \eqref{eq:jacobi} is spectral only if $G$ is analytic. By \eqref{eq:G} it is
not: $G$ carries fractional powers $r^{p\nu}$. For a \emph{single} mode
($c_k=0$, $k\ge2$) we have $G(r)=b_1^2A_{\nu_1}(\lambda r^2)^2$, analytic in
$r^2$, and \eqref{eq:jacobi} is spectral --- but a single mode is not the
situation of interest. For a general solution, convergence degrades to algebraic.

\begin{proposition}[Rationalizing substitution]\label{prop:sub}
Under $t=r^{\nu}$, equivalently $r=t^{1/\nu}$, the edge integral becomes
\begin{equation}\label{eq:sub}
  \int_0^R r^{2\nu-2}G(r)\dd r
  \;=\;\frac1\nu\int_0^{R^\nu} t^{\,1-1/\nu}\,\widetilde G(t)\dd t,
  \qquad
  \widetilde G(t)=\sum_{p,q}g_{pq}\,t^{\,p+2q/\nu},
\end{equation}
a Jacobi-weight integral with exponent $\beta'=1-1/\nu$. The corner family
$r^{p\nu}$ maps to the integer powers $t^{p}$.
\end{proposition}
\begin{proof}
$\dd r=\frac1\nu t^{1/\nu-1}\dd t$ and $r^{2\nu-2}=t^{2-2/\nu}$; the exponents
add to $(2-2/\nu)+(1/\nu-1)=1-1/\nu$. Substituting $r^{p\nu+2q}=t^{p+2q/\nu}$
into \eqref{eq:G} gives $\widetilde G$.
\end{proof}

\noindent
Admissibility $\beta'>-1$ again requires exactly $\nu>1/2$.

\paragraph{The caveat.} The substitution rationalizes the corner family
completely, but maps the Bessel family $r^{2q}\mapsto t^{2q/\nu}$, which is
integral only when $2/\nu\in\mathbb{Z}$ --- equivalently $\alpha=m\pi/2$ for an
integer $m$. The $270^\circ$ corner is such a case:
\begin{equation}
  \alpha=\tfrac32\pi \;\Longrightarrow\; \nu=\tfrac23 \;\Longrightarrow\;
  2/\nu=3\in\mathbb{Z}, \qquad \widetilde G(t)=\sum g_{pq}t^{\,p+3q}
  \ \ \text{(a power series in $t$).}
\end{equation}
For generic $\alpha$ the residual exponents $2q/\nu$ are fractional, but they are
\emph{high order}: the leading one is $2/\nu>2$, so $\widetilde G\in C^{2}$ at
worst and Gauss convergence is algebraic of high order rather than spectral.
This is confirmed in Section~\ref{sec:num}(c): the machine precision at $n=8$
is specific to $\alpha=\frac32\pi$, while generic $\alpha$ converges at roughly
two digits per node doubling --- still reaching $10^{-12}$ by $n=32$, which is
cheap enough that the distinction is of no practical consequence.

\section{Sensitivity to $\nu$}\label{sec:sens}

The weight exponent must use $\nu$ at full precision. Perturbing $\nu\to\nu+\delta$
leaves a residual factor $r^{-2\delta}$ in the integrand, which is unbounded as
$r\to0$ for $\delta>0$ and destroys the analyticity the rule depends on.
Measured at $\alpha=1.5\pi$ (true $\nu=2/3$), mode $(1,1)$:
\medskip

\begin{center}
\begin{tabular}{lcc}
\toprule
$\nu$ used in $\beta$ & 28 nodes & 192 nodes \\
\midrule
$0.6667$ (exact)      & $-7.5\times10^{-15}$ & $-3.1\times10^{-15}$\\
$0.667$               & $6.6\times10^{-5}$   & $1.7\times10^{-5}$\\
$0.600$               & $-1.8\times10^{-2}$  & $-4.7\times10^{-3}$\\
\bottomrule
\end{tabular}
\end{center}
\medskip

\noindent
A relative error of $3\times10^{-4}$ in $\nu$ costs four digits and reduces
convergence to algebraic. Since $\nu=\pi/\alpha$ is exact from the geometry this
costs nothing --- but it rules out tabulated, rounded, or margin-padded exponents
of the kind used for grading orders.

\section{Numerical verification}\label{sec:num}

All errors are relative, against closed-form sector eigenfunctions of unit norm
(exact $L^2$ norms, so the reference is analytic rather than quadrature-based).

\paragraph{(a) Single mode, $270^\circ$ corner, $x_0$ at the failing default.}
Cauchy data treated as a black box; $2n_{\text{edge}}+n_{\text{arc}}$ nodes total.
\medskip

\begin{center}
\begin{tabular}{lcc}
\toprule
nodes & Gauss--Jacobi & existing Kress rule\\
\midrule
16   & $-7.1\times10^{-9}$  & ---\\
28   & $-7.5\times10^{-15}$ & ---\\
2264 & ---                  & $2.2\times10^{-7}$ (plateau)\\
\bottomrule
\end{tabular}
\end{center}

\paragraph{(b) Multi-exponent corner data.} A five-term expansion
\eqref{eq:expansion} times an analytic tail, integrated against an mpmath
reference at 40 digits, at $\nu=2/3$:
\medskip

\begin{center}
\begin{tabular}{lcc}
\toprule
$n$ & Jacobi in $r$ & Jacobi in $t=r^{\nu}$ \\
\midrule
8  & $-4.6\times10^{-3}$ & $4.0\times10^{-15}$\\
16 & $-1.2\times10^{-3}$ & $7.3\times10^{-15}$\\
64 & $-7.2\times10^{-5}$ & $1.6\times10^{-13}$\\
\bottomrule
\end{tabular}
\end{center}
\medskip

\noindent
Plain Jacobi in $r$ is algebraic, as Section~\ref{sec:sub} predicts; the
substituted rule is spectral and \emph{degrades} slowly past $n\approx32$ from
roundoff amplification, so small $n$ is both cheaper and more accurate.

\paragraph{(c) Generic $\alpha$: is $\nu=2/3$ special?} The same five-term data,
substituted rule, against an \emph{exact} reference (see the warning below).
\medskip

\begin{center}
\begin{tabular}{lccccc}
\toprule
$\alpha$ & $2/\nu$ & $n=8$ & $n=16$ & $n=32$ & $n=64$\\
\midrule
$1.5\pi$   & $3.000$ (int.) & $2.0\times10^{-15}$ & $2.0\times10^{-15}$ & $-2.6\times10^{-14}$ & $1.7\times10^{-13}$\\
$1.33\pi$  & $2.667$ & $2.8\times10^{-8}$ & $2.7\times10^{-10}$ & $2.7\times10^{-12}$ & $5.1\times10^{-14}$\\
$1.25\pi$  & $2.500$ & $4.7\times10^{-8}$ & $5.1\times10^{-10}$ & $5.9\times10^{-12}$ & $5.0\times10^{-14}$\\
$1.6\pi$   & $3.200$ & $-5.6\times10^{-9}$ & $-3.4\times10^{-11}$ & $-1.6\times10^{-13}$ & $-7.7\times10^{-13}$\\
$1.7\pi$   & $3.400$ & $-5.2\times10^{-9}$ & $-2.7\times10^{-11}$ & $-3.4\times10^{-13}$ & $-9.8\times10^{-13}$\\
$1.85\pi$  & $3.700$ & $-1.5\times10^{-9}$ & $-6.0\times10^{-12}$ & $1.9\times10^{-14}$ & $-2.5\times10^{-12}$\\
$1.95\pi$  & $3.900$ & $-1.4\times10^{-10}$ & $-2.0\times10^{-13}$ & $-4.0\times10^{-13}$ & $-1.2\times10^{-11}$\\
\bottomrule
\end{tabular}
\end{center}
\medskip

\noindent
Only $\alpha=\frac32\pi$ is machine-exact at $n=8$, confirming
Section~\ref{sec:sub}: the integer $2/\nu$ is genuinely special. Every other
angle converges at $\approx$ two digits per doubling --- high-order algebraic,
exactly as predicted --- and all reach $10^{-12}$ or better by $n=32$. The mild
degradation past $n\approx32$ is roundoff, not truncation.

\paragraph{Warning: mpmath's \texttt{quad} is not a safe reference here.}
The endpoint behaviour $r^{2\nu-2}\to r^{-1}$ as $\nu\to\frac12$ defeats
\texttt{mp.quad} even at 40 digits. Measured against the exact termwise value
$\int_0^1 r^a\dd r=(a+1)^{-1}$, its error reaches $-4.7\times10^{-5}$ at
$\alpha=1.85\pi$ and $-4.3\times10^{-2}$ at $\alpha=1.95\pi$. Using it as ground
truth produces convincing but entirely spurious ``plateaus'' in the rule under
test. Any reference for this integral should be closed-form.

\paragraph{A trap worth recording.} At $\alpha=3\pi/2$ the two edges carry
$r_N=\operatorname{Im}x_0$ and $-\operatorname{Re}x_0$ against identical
$(\partial_n u)^2$, so any $x_0$ on the diagonal $\operatorname{Re}x_0=
\operatorname{Im}x_0$ cancels the singular contribution \emph{identically}. A
test using such an $x_0$ silently measures the smooth part of the boundary only.
Corner-quadrature experiments should always report the edge/arc split.

\section{Recommended construction}

For each corner with $\nu<1$, and each straight edge meeting it:
\begin{enumerate}
\item Work in corner-local coordinates $(r,\theta)$; never form
$p_0+\tau L\hat d$.
\item Panel the edge, substitute $t=r^{\nu}$ with $\nu=\pi/\alpha$ exact.
\item Apply Gauss--Jacobi with $\beta'=1-1/\nu$ and $n=8$--$16$ nodes.
\end{enumerate}
Smooth portions of $\partial\Om$ retain plain Gauss--Legendre. Edges meeting two
singular corners are split, one panel per endpoint.

This removes the requirement that $x_0$ sit at a corner, and with it the
obstruction to multi-corner domains: no partition of unity over corners is
needed. The requirement $\nu>1/2$ is not restrictive --- it is exactly the
condition for \eqref{eq:edge} to converge at all --- but the slit limit
$\nu\to1/2$ remains delicate, both weights degenerating as $\beta,\beta'\to-1$.

\paragraph{Not yet verified.} A real \texttt{FourierBesselBasis} through this
rule (as opposed to closed-form or synthetic data), and a genuinely multi-corner
domain.


\section{Addendum: what implementation changed}\label{sec:addendum}

The analysis above was written before the rule was built. Its core result held:
a corner-adapted rule reaches $10^{-14}$ where the graded rule plateaus at
$10^{-7}$, with fewer nodes. What follows corrects the parts that did not survive
contact with a real domain, and records what could only be learned by measuring.
Everything here is implemented in \verb|lappy.quad| and
\verb|lappy.eigfun_integrals|, and tested in \verb|tests/test_eigfun_integrals.py|.

\subsection{The substitution is $t=r^{1/q}$, and $t=r^{\nu}$ is a trap for curved edges}

Section~\ref{sec:sub} identified the exponent family as $\{k\nu+2q\}$ and
concluded that $t=r^{\nu}$ rationalizes it when $2/\nu\in\mathbb{Z}$ --- for
reentrant $\alpha$, only $\alpha=3\pi/2$. That is correct \emph{for a straight
edge}, and it is the wrong family for a curved one.

On a curved edge two things change, both adding \emph{integer} powers of
arclength $s$. First, $r_N$ is no longer constant: with $x_0$ at the corner it is
$(\kappa_0/2)s^2+O(s^3)$ --- vanishing to second order rather than identically,
so Proposition~\ref{prop:rn}'s escape degrades gracefully rather than failing.
Second, $r=|x(s)-\text{corner}|=s-(\kappa_0^2/24)s^3+\cdots$, so $r$ is not
arclength. The family becomes $\{k\nu+m\}$, $m\in\mathbb{Z}_{\ge0}$, including
\emph{odd} powers. Under $t=r^{\nu}$ that maps to $\{k+m/\nu\}$ whose leading
residual $t^{1/\nu}$ has $1/\nu\in(1,2)$: only $C^1$. Measured, it never reaches
$10^{-13}$ at any order up to 64. Under $t=r^{1/q}$ with $\nu=p/q$ in lowest
terms it maps to $\{kp+mq\}$ --- all integers --- and is exact by order 6--14.

\emph{But $t=r^{1/q}$ is unusable as written.} Since $\tau=t^{q}$, the innermost
node behaves as $\tau_{\min}\sim t_{\min}^{q}$:

\begin{center}
\begin{tabular}{lcl}
\toprule
$\alpha$ & $q$ & $\tau_{\min}$ at order 24\\
\midrule
$3\pi/2$  & 3  & $1.4\times10^{-8}$\\
$5\pi/4$  & 5  & $1.1\times10^{-10}$ \quad(below the floor)\\
$7\pi/4$  & 7  & $4.7\times10^{-19}$\\
$11\pi/6$ & 11 & $1.6\times10^{-29}$\\
\bottomrule
\end{tabular}
\end{center}

\noindent
Those nodes round onto the corner once mapped through a segment's
parametrization --- the very defect Remark~1 identifies. The rule is exact as an
abstract statement about $[0,1]$ and worthless on a boundary.

\subsection{What works: exactness from the weights, node placement from $\nu$}

The resolution separates the two roles the substitution was doing. Nodes are
placed by $t=r^{\nu}$, whose crowding is mild ($\tau_{\min}$ stays at
$10^{-5}$--$10^{-7}$). Exactness comes from the \emph{weights}: the exponent set
$\{\gamma+j\nu+m\}$ is known, and every moment is closed-form,
$\int_0^1 t^{e}\dd t=(e+1)^{-1}$, so the weights are obtained by solving a linear
system. This also covers irrational $\nu$, which the substitution cannot touch at
all --- and irrational $\nu$ is the \emph{generic} case for a curved corner, whose
angle is fixed by the geometry rather than chosen (two overlapping circles give
$\nu=0.657360$).

One trap: solving the square system (as many exponents as nodes) is exact on its
span but has $\mathrm{cond}\sim10^{13}$ at order 12 and $10^{19}$ at order 16,
with weights reaching $-10^{4}$. Since $\sum|w|$ multiplies every roundoff error
in the integrand, that imports a $10^{-12}$ floor. Using \emph{fewer exponents
than nodes} --- a minimum-norm least-squares solve --- keeps $\sum|w|=1.0$ with
exactness intact.

\subsection{The edge type matters, and conflating them costs seven orders}

A straight edge admits only \emph{even} integer powers ($r_N$ exactly constant,
$r$ exactly arclength), so $\{\gamma+j\nu+2q\}$; odd powers are curved-only.
Scored against the curved set, the substitution rule at $\alpha=3\pi/2$ reads
$10^{-8}$ --- against a measured $2.4\times10^{-15}$ on a real sector. Any error
indicator must use the family belonging to the edge it is scoring.

\subsection{Panel length must be capped by clearance, not by a fraction}

Section~6's recipe panels the edge but says nothing about how long a panel may
be. The corner expansion is valid only within the largest disk about the corner
inside $\Om$; a panel reaching past it is under-resolved, because the rule
clusters its nodes \emph{at} the corner and cannot resolve the $\sqrt\lambda$
oscillation over the remaining arclength. Measured on a $1\times N$ polyomino
strip with one cell below an end (so the corner's edge has length $N-1$ against a
clearance of 1), against its exact eigenfunction:

\begin{center}
\begin{tabular}{lcccc}
\toprule
$N$ & no cap & $0.9\times$ clearance\\
\midrule
2  & $6.5\times10^{-14}$ & $5.3\times10^{-15}$\\
8  & $2.3\times10^{-2}$  & $1.6\times10^{-15}$\\
16 & $7.7\times10^{-2}$  & $1.3\times10^{-15}$\\
24 & $7.4\times10^{-3}$  & $4.4\times10^{-16}$\\
\bottomrule
\end{tabular}
\end{center}

\noindent
Twelve orders. A fixed \emph{fraction} of the edge cannot substitute --- being a
fraction of the edge, it grows with the edge and stays too long ($0.25$ of the
edge still gives $2.6\times10^{-4}$ at $N=16$). The cap must be geometric.

\subsection{Two properties that look like bugs and are not}

$\sum_i w_i=1$ holds only when the constant function lies in the rule's exact
class. Renormalizing would destroy exactness on the singular class the rule
exists for; $\sum_i w_i-1$ is instead its cheapest self-diagnostic. And accuracy
is \emph{not monotone} in order: past a $\nu$-dependent threshold the integrand's
dynamic range ($\tau^{2\nu-2}\sim6\times10^{8}$ at $\tau\sim10^{-9}$ against a
correspondingly tiny weight) makes roundoff dominate. At $\nu=0.526$ the best
order is 16 while the floor-based cap is 75, and using the cap gives
$4.2\times10^{-9}$. Order selection must therefore scan a calibrated curve, not
increase the order until a target is met.

\subsection{Scope, sharpened}

The rule is matched to \emph{eigenfunction}-level Cauchy data, whose local
structure is purely the corner family. It is \emph{not} appropriate for a
basis-level $N\times N$ Gram matrix: columns centred at other corners are plain
analytic at this one, with $O(1)$ amplitude at every corner simultaneously.
Measured, the corner rule beats the graded rule by 3--8 orders on corner-family
blocks and loses by 2--4 on mixed ones. That is why \verb|lappy| retired the
basis-level path rather than adapting it.

Finally, a warning about instruments. \verb|mpmath.quad| is not a valid reference
for these integrands: as $\nu\to1/2$ the endpoint behaviour approaches $r^{-1}$
and it errs by $4\times10^{-2}$ even at 40 digits, manufacturing convincing but
entirely spurious convergence plateaus. Every reference used in the tests is
closed-form.

\end{document}
